\scanpage{61}
% archive sheet number: 60
finies existent, et que $T_X$ est exact à g., on est
ramené à vérifier que tt couple $(P,\unsure{\psi|u|\varphi}\in T_X(P))$
est \deleted{\unsure{majoré}} dominé par un ``couple minimal''.

Interprétons en effet le couple en question par
$\varphi:X\times_S X\longrightarrow G$ satisfaisant $\varphi(x,x)=e$, alors \deleted{\unsure{$G$}} les sous-groupes fermés $G_\alpha$ de $G$ tels que $\varphi$ se factorise par $G_\alpha\times_S G_\alpha$
\deleted{forment} un ens filtrant décroissant, et \unsure{because le noeth|because la noeth}
il y a un objet minimal, soit $G'$. Le couple
$(G',\varphi')$ correspondant est le minorant cherché.

\deleted{Montrons que le \unsure{groupe} pro-objet \inserted{\unsure{strict}} $(G_\alpha,\varphi_\alpha)$ qui représente $T_X$ est \unsure{ess. constant}.}
Nous allons
\deleted{montrer que} \inserted{mettre} le foncteur précédent \deleted{\unsure{est}} sous une
forme simple lorsque on y astreint $G$ à être un
schém. \uline{abélien} $A$, donc $A=\underline{\mathrm{Pic}}^0_{B/S}$, où $B=\widehat A=\underline{\mathrm{Pic}}^0_{\unsure{B|A}/S}$.
En effet, \deleted{alors} \deleted{soit $T_X(B)$ \unsure{l'alg de} \unsure{donne la} \unsure{fibre} \unsure{point}}
prendre pour simplifier $S'=X$, et faire le chgt de base \unsure{$S'\to S$|$S'$} 
\[
\begin{aligned}
\unsure{\pi_X|\Pi_X|T_X}(A)&\simeq\mathrm{Hom}'_S(X\times_S X,A)\simeq\mathrm{Hom}_{S'\text{-}\unsure{\mathrm{point\acute es}|\mathrm{point\acute e}}}(X_{S'},A_{S'})=\mathrm{Hom}_{S'\text{-}\mathrm{point}}(X_{S'},\underline{\mathrm{Pic}}^0_{B_{S'}/S'})\\
&\simeq\underbrace{\unsure{\mathscr B|\mathcal B}_{S'}(X_{S'},B_{S'})}_{\substack{\text{\uline{classes de corr.}}\\\text{\uline{divisorielles}}}}\simeq\mathrm{Hom}_{S'\text{-}\mathrm{point.}}(B_{S'},\underline{\mathrm{Pic}}_{X_{S'}/S'})=\underline{\mathrm{Hom}}_{S'\text{-}\unsure{\mathrm{gr}|\mathrm{point}}}(B_{S'},(\underline{\mathrm{Pic}}_{X/S})_{S'})\\
&\underset{\substack{\text{d'après le}\\\text{théorème de}\\\text{la rigidité}}}{\simeq}\mathrm{Hom}_{S\text{-groupes}}(B,\underline{\mathrm{Pic}}_{X/S})\qquad\text{donc}
\end{aligned}
\]
\[
\boxed{\pi_X(A)\simeq\mathrm{Hom}_{S\text{-gr}}(\widehat A,\underline{\mathrm{Pic}}_{X/S})}
\]
qui est un résultat du type de ``dualité de Cartier''.

Or supposons que $\underline{\mathrm{Pic}}_{X/S}$ contienne un sous-schéma
\deleted{\unsure{$B_{X/S}$}} \inserted{\unsure{disons|de s}} \inserted{l'un ss-groupe \unsure{sch.}} \uline{abélien} $\underline{\mathrm{Pic}}^0_{X/S}$, \unsure{Alors|dans} \unsure{utilisant|satisfaisant} \deleted{les} \deleted{\unsure{ss}}
``pts d'ordre fini''
